% formal.tex — page furniture for formal business documents.
%
%   pandoc doc.md -H formal.tex --pdf-engine=xelatex -o doc.pdf
%
% Provides a footer imprint, a signature line, and an optional letterhead.
% Everything here is EMPTY BY DEFAULT and opt-in: including this file does not
% stamp anyone's details onto your document. Set the values yourself, either in
% your own header file or in the document's metadata.

\usepackage[a4paper,top=30mm,bottom=25mm,left=25mm,right=20mm,footskip=15mm]{geometry}
\usepackage{fancyhdr}
\usepackage{xcolor}
\usepackage{graphicx}

% — Colours ————————————————————————————————————————————————
% Neutral greys. Override in your own header to match your house style.
\providecolor{doc-ink}{HTML}{1A1A1A}       % body text
\providecolor{doc-secondary}{HTML}{666666} % footer, secondary text
\providecolor{doc-hairline}{HTML}{D8D8D8}  % rules and dividers

% — Imprint ————————————————————————————————————————————————
% Many jurisdictions require certain particulars on business correspondence.
% German law, for instance, requires a GmbH's Geschäftsbriefe to carry the
% firm, legal form, seat, register court, register number and all managing
% directors (§ 35a Abs. 1 S. 1 GmbHG); comparable duties exist elsewhere.
%
% This is a TYPESETTING component, not legal advice, and it does not know what
% your jurisdiction requires of you. Redefine \docimprint with your own text:
%
%   \renewcommand{\docimprint}{%
%     \footnotesize\color{doc-secondary}%
%     \begin{tabular}[b]{@{}l@{}}
%     Example Ltd · 1 Example Street · 12345 Example City\\
%     Directors: A. Example · Reg. 12345, Example Registry
%     \end{tabular}}
%
% Left empty, the footer simply carries a page number.
\providecommand{\docimprint}{}

\makeatletter
\renewcommand{\footrule}{{\color{doc-hairline}%
  \vskip-\footruleskip\vskip-\footrulewidth
  \hrule\@width\headwidth\@height\footrulewidth
  \vskip\footruleskip}}
\makeatother

\pagestyle{fancy}
\fancyhf{}
\renewcommand{\headrulewidth}{0pt}
\renewcommand{\footrulewidth}{0.4pt}
\fancyfoot[L]{\docimprint}
\fancyfoot[R]{\footnotesize\color{doc-secondary}\thepage}

% The `plain` page style applies to pages carrying \maketitle. Without this
% mirror, page one — precisely the page a letter's particulars belong on —
% would silently lose the imprint while every later page kept it. That
% asymmetry is easy to ship without noticing, because the document looks fine
% unless you check page one specifically.
\fancypagestyle{plain}{%
  \fancyhf{}
  \renewcommand{\headrulewidth}{0pt}
  \renewcommand{\footrulewidth}{0.4pt}
  \fancyfoot[L]{\docimprint}
  \fancyfoot[R]{\footnotesize\color{doc-secondary}\thepage}}

\usepackage{setspace}
\setstretch{1.15}

% — Letterhead (opt-in) ————————————————————————————————————
% Define \docletterhead with your own logo, then call it in the document.
% Not set automatically: a resolution or minute usually wants no letterhead.
%
%   \renewcommand{\docletterhead}{%
%     \noindent\includegraphics[height=9mm]{logo.pdf}\par\vspace{8mm}}
\providecommand{\docletterhead}{}

% — Signature line ————————————————————————————————————————
% \signatureline{place and date}{name and capacity}
%
% The minipage is the point. Without it the block is breakable, and a signature
% block that breaks is a real defect rather than a cosmetic one: the date line
% ends up at the foot of one page while the rule you actually sign and the name
% beneath it sit alone on the next — a signature detached from the document it
% attests to, and from its own date.
%
% This is not hypothetical — it is why the minipage is here. Note that no
% text-layer check can see it: nothing is missing, it is merely in the wrong
% place. The way to catch it is to ask, page by page, whether the date line and
% the name ended up on the same page. tests/run.sh does exactly that.
%
% A minipage cannot break across pages. Either the whole block fits where it is,
% or the whole block moves to the next page. The leading space stays OUTSIDE it,
% so that it collapses at a page break instead of pushing the block down the new
% page.
\providecommand{\signatureline}[2]{%
  \par\vspace{10mm}%
  \noindent\begin{minipage}{\linewidth}%
    \noindent{\small #1}\par
    \vspace{16mm}\noindent{\color{doc-ink}\rule{70mm}{0.4pt}}\\
    {\small #2}%
  \end{minipage}}
